\item \textbf{Problema de repaso: Vibraciones en la cuerda de un yo-yo.}
El extremo superior de una cuerda de un yo-yo se mantiene en reposo. De hecho, el yo-yo es más pesado que la cuerda. Parte del reposo y se mueve hacia abajo con aceleración constante de $0.800\text{ m/s}^2$ mientras se desenrolla la cuerda. El roce de la cuerda contra el borde del yo-yo excita vibraciones de onda estacionaria transversales en la cuerda. Ambos extremos de la cuerda son nodos, incluso cuando la longitud de la cuerda aumenta. Considere el instante $1.20\text{ s}$ después de que comienza el movimiento a partir del reposo.
\begin{enumerate}
    \item Demuestre que la razón de cambio con el tiempo de la longitud de onda del modo fundamental de oscilación es $1.92\text{ m/s}$.
    \item \textbf{¿Qué pasaría si?} ¿La razón de cambio de la longitud de onda del segundo armónico también es $1.92\text{ m/s}$ en este momento? Explique.
    \item \textbf{¿Qué pasaría si?} El experimento se repite después de agregar más masa al cuerpo del yo-yo. La distribución de masa se mantiene igual, de modo que el yo-yo todavía se mueve con aceleración hacia abajo de $0.800\text{ m/s}^2$. En el punto $1.20\text{ s}$, ¿la razón de cambio de la longitud de onda fundamental de la cuerda en vibración aún es igual a $1.92\text{ m/s}$? Explique.
    \item ¿La razón de cambio de la longitud de onda del segundo armónico es la misma que en el inciso (2)? Explique.
\end{enumerate}